% ======================================================
%  Rozdział — Zagadki Logiczne
%
%  GENERATED by tools/convert_deck.py from areas/3-zagadki-logiczne.tex.
%  Edit the deck source and re-run, or convert this chapter to hand-maintained
%  and delete it from the script's AREA_META -- but do not edit both.
%
%  The `% -- NN --` comments are the deck's own global exercise numbers, kept
%  so a reader of either source can find the same exercise in the other. The
%  book itself numbers frames per chapter: within a chapter, frame N is
%  exercise N, because every frame carries exactly one exercise except the last.
% ======================================================

\chapter{Zagadki Logiczne}

Sylogizmy, pułapki słowne i dedukcja. Tu liczy się to, co przeczytasz, nie to, co policzysz.

% -- 19 --
\begin{fr}
  \exercise[\CategoryBadge[LogicColor!20]{Sylogizm}]{1}{2 min}{%
      Wszystkie Bloopy są Razzies.\par
      Wszystkie Razzies są Lazzies.\par
      Czy wszystkie Bloopy są Lazzies?}
\end{fr}

% -- 20 --
\begin{fr}
  \ans{Tak!}{Bloopy $\to$ Razzies $\to$ Lazzies.}
  \exercise[\CategoryBadge[LogicColor!20]{Pu\l{}apka}]{2}{2 min}{%
      Rolnik ma 17 owiec.\par
      Wszystkie poza 9 zdychaj\k{a}.\par
      Ile owiec zostaje?}
\end{fr}

% -- 21 --
\begin{fr}
  \ans{9}{\emph{Poza 9} znaczy: 9 prze\. zy\l{}o.}
  \exercise[\CategoryBadge[LogicColor!20]{Pu\l{}apka}]{2}{3 min}{%
      Dwie monety daj\k{a} razem 30 groszy.\par
      Jedna NIE jest 10-groszowk\k{a}.\par
      Jakie to monety?}
\end{fr}

% -- 22 --
\begin{fr}
  \ans{20~gr~+~10~gr}{Ta \emph{jedna} (20-gr) naprawde nie jest 10-gr.}
  \exercise[\CategoryBadge[LogicColor!20]{Dedukcja}]{2}{3 min}{%
      Adam, Bartek, Celina: pi\l{}ka, p\l{}ywanie, bieganie.\par
      Adam nie lubi pi\l{}ki. Bartek p\l{}ywa.\par
      Co lubi Celina?}
\end{fr}

% -- 23 --
\begin{fr}
  \ans{Pi\l{}k\k{e} no\. zn\k{a}}{Bartek\,=\,p\l{}ywanie; Adam\,=\,bieganie; Celina\,=\,pi\l{}ka.}
  \exercise[\CategoryBadge[LogicColor!20]{Geometria}]{3}{5 min}{%
      Zegar pokazuje 3:15.\par
      Jaki k\k{a}t tworz\k{a} wskaz\'{o}wki?}
\end{fr}

% -- 24 --
\begin{fr}
  \ans{$7{,}5^\circ$}{Godzinowa: $97{,}5^\circ$ · minutowa: $90^\circ$ $\to$ r\'{o}\. znica $7{,}5^\circ$.}
  \exercise[\CategoryBadge[LogicColor!20]{Pu\l{}apka}]{3}{5 min}{%
      5 maszyn robi 5 zabawek w 5 minut.\par
      Ile czasu zajmie 100 maszynom\par
      zrobi\'c 100 zabawek?}
\end{fr}

% -- 25 --
\begin{fr}
  \ans{5 minut}{Ka\. zda maszyna niezale\. znie --- 100 r\'{o}wnolegle.}
  \exercise[\CategoryBadge[LogicColor!20]{Pu\l{}apka}]{3}{5 min}{%
      W wy\'scigu wyprzedzasz osob\k{e}\par
      na 2.~miejscu.\par
      Na kt\'{o}rym miejscu jeste\'s?}
\end{fr}

% The chapter's last frame carries the last answer alone.
\begin{fr}
  \ans{2.~miejsce}{Zajmujesz jej pozycj\k{e} --- nie wyprzedzasz lidera.}
\end{fr}
